\section{Three agreements}\label{sec:agreements}

A protocol that pays for use must bring its participants to agree on three different things. They differ in kind
and rest on different assumptions, and treating them as one problem hides where the difficulty lies.
Table~\ref{tab:agreements} separates them.

\begin{table}[!htb]
\centering\small
\begin{tabularx}{\textwidth}{@{}p{0.23\textwidth}XX@{}}
\toprule
\textbf{What is agreed} & \textbf{How} & \textbf{Assumption} \\
\midrule
Which constructions exist and who registered each first &
Signed, hashed registration, ordered by the host chain &
The host chain's own, for the order. Who said what needs none: unforgeable signatures settle it for any number of
traitors \citep{lamport1982byzantine} \\
\addlinespace
Whether a use happened &
Re-execution of the receipt against the committed text; paid spot checks; a challenge on mismatch &
An honest node is among those that examine the receipt \citep{canetti2011delegation}; checking is paid for
\citep{luu2015verifier}; a disputed finding needs honest voting stake \\
\addlinespace
Whether the Bonfires that hold a construction are independent &
Cannot be re-executed. Anchored in stake that can be lost, and in a reward that a small group does not reach &
Economic: posing as many costs more than it earns \\
\bottomrule
\end{tabularx}
\caption{The three agreements of the protocol, each with the assumption it rests on.}\label{tab:agreements}
\end{table}

The first agreement is on the order of signed messages, which existing ledgers already provide
\citep{nakamoto2008bitcoin}; the protocol takes it from a host chain. The second is the verification of a computation over committed inputs that the protocol requires to be
deterministic, for which a single honest party that re-executes is enough to find
an error \citep{canetti2011delegation}, though a finding that the accused disputes is settled by a vote. Where
checking is costly, rational parties have an incentive to accept results unchecked \citep{luu2015verifier}, so the
protocol pays for spot checks whether or not they find a fault.

The third is where the difficulty lies. That a use happened does not show that it was anyone's but the
claimant's. Any party can produce correct and pointless conversations in any number, and no computation tells them
from useful ones. So the protocol does not count uses. It counts the independent Bonfires that have taken a
construction up and still hold it (Section~\ref{sec:reward}), and it pays nothing until they are several. One party
can pose as many Bonfires, and whether it has done so cannot be computed. That question is answered by stake: each
Bonfire that counts is bonded, what can be earned is capped by the stake committed, and stake is lost when several
Bonfires are found to be one party (Section~\ref{sec:bond}). The first two agreements are borrowed; the protocol
adds only that structuring discourse into constructions is what makes the second available.

\begin{figure}[!htb]
\centering
% fig:life -- the life of a construction
\begin{tikzpicture}[
  font=\small,
  stage/.style={draw=cxnnavy, fill=cxngreen, rounded corners=2pt, minimum width=2.05cm, minimum height=0.8cm,
               align=center},
  arr/.style={-{Stealth[length=5pt]}, thick, cxnnavy},
  note/.style={font=\scriptsize, text=black!70, align=center}
]
  \node[stage] (pub) {1. Make};
  \node[stage, right=0.45cm of pub] (ado) {2. Serve};
  \node[stage, right=0.45cm of ado] (enc) {3. Look up};
  \node[stage, right=0.45cm of enc] (chk) {4. Adopt};
  \node[stage, right=0.45cm of chk] (rep) {5. Count};
  \node[stage, right=0.45cm of rep] (set) {6. Pay};
  \draw[arr] (pub) -- (ado);
  \draw[arr] (ado) -- (enc);
  \draw[arr] (enc) -- (chk);
  \draw[arr] (chk) -- (rep);
  \draw[arr] (rep) -- (set);
  \node[note, below=0.12cm of pub] {registered,\\published};
  \node[note, below=0.12cm of ado] {nodes\\pledge};
  \node[note, below=0.12cm of enc] {from\\a gap};
  \node[note, below=0.12cm of chk] {trial, then\\confirmation};
  \node[note, below=0.12cm of rep] {Bonfires that\\hold it};
  \node[note, below=0.12cm of set] {epoch budget\\divided};
  \draw[arr, modelred] (chk.north) -- ++(0,0.55) -| node[pos=0.25, above, font=\scriptsize, text=modelred]
    {a variant, registered with its parent} (pub.north);
\end{tikzpicture}

\caption{The life of a construction. Members' discourse yields it and their Bonfire registers and publishes it;
nodes pledge to the Bonfire and serve it; another Bonfire's gap becomes a lookup; the construction is tried and
adopted; each epoch the Bonfires that hold it are counted; the epoch budget is divided. A Bonfire that changes it
registers a variant with its parent.}\label{fig:life}
\end{figure}
